% =====================================================================
%  quizgen — Implementation Report
%  Automated Test Assembly: selecting exams from a question bank
%
%  Compile with:  pdflatex implementation.tex   (run twice for refs)
%             or:  tectonic implementation.tex
% =====================================================================
\documentclass[11pt,a4paper]{article}

\usepackage[utf8]{inputenc}
\usepackage[T1]{fontenc}
\usepackage[margin=2.4cm]{geometry}
\usepackage{amsmath,amssymb}
\usepackage{booktabs}
\usepackage{longtable}
\usepackage{array}
\usepackage{enumitem}
\usepackage{xcolor}
\usepackage{listings}
\usepackage{tikz}
\usepackage{caption}
\usepackage[hidelinks]{hyperref}
\usepackage{microtype}

\usetikzlibrary{arrows.meta,positioning,fit,backgrounds,calc,shapes.geometric}

% ---- colours ---------------------------------------------------------
\definecolor{cstage}{RGB}{222,235,247}
\definecolor{cstageb}{RGB}{ 49,130,189}
\definecolor{cdata}{RGB}{229,245,224}
\definecolor{cdatab}{RGB}{ 49,163, 84}
\definecolor{cctrl}{RGB}{253,224,221}
\definecolor{cctrlb}{RGB}{222, 45, 38}
\definecolor{cgrey}{RGB}{99,99,99}
\definecolor{ckw}{RGB}{0,90,160}
\definecolor{cstr}{RGB}{160,60,0}
\definecolor{ccmt}{RGB}{100,120,100}

% ---- listings (python / yaml) ---------------------------------------
\lstdefinestyle{py}{
  language=Python,
  basicstyle=\ttfamily\footnotesize,
  keywordstyle=\color{ckw}\bfseries,
  stringstyle=\color{cstr},
  commentstyle=\color{ccmt}\itshape,
  numberstyle=\tiny\color{cgrey},
  showstringspaces=false,
  breaklines=true,
  frame=single,
  rulecolor=\color{cgrey!40},
  framesep=4pt,
  columns=fullflexible,
  keepspaces=true,
}
\lstset{style=py}

\newcommand{\code}[1]{\texttt{\small #1}}
\newcommand{\mod}[1]{\texttt{\small #1}}

% ---- TikZ node styles ------------------------------------------------
\tikzset{
  stage/.style={rectangle, rounded corners=3pt, draw=cstageb, fill=cstage,
                line width=0.8pt, align=center, inner sep=6pt, font=\small\bfseries,
                minimum height=10mm, minimum width=26mm},
  data/.style={rectangle, draw=cdatab, fill=cdata, line width=0.7pt,
               align=center, inner sep=5pt, font=\footnotesize,
               minimum height=8mm},
  ctrl/.style={rectangle, rounded corners=2pt, draw=cctrlb, fill=cctrl,
               line width=0.7pt, align=center, inner sep=4pt, font=\footnotesize},
  arr/.style={-{Stealth[length=2.4mm]}, line width=0.8pt, draw=cgrey},
  darr/.style={-{Stealth[length=2.2mm]}, line width=0.7pt, draw=cdatab, densely dashed},
}

% =====================================================================
\begin{document}

\begin{titlepage}
\centering
\vspace*{2.2cm}
{\Huge\bfseries quizgen\par}
\vspace{6mm}
{\LARGE Automated Test Assembly from a Question Bank\par}
\vspace{3mm}
{\large Selecting blueprint-compliant exams from a pool of\\
tagged candidate questions via optimisation\par}
\vspace{1.4cm}
{\large\bfseries Implementation Report\par}
\vspace{4mm}
{\normalsize Master's Thesis --- KWARC Group\\
Friedrich-Alexander-Universit\"at Erlangen-N\"urnberg\par}
\vfill
{\small This report documents the implemented logic of the \code{quizgen}
codebase --- a selection-only tool that turns an existing question bank into
exam papers --- the algorithms behind each stage, and the literature each
technique draws on.\par}
\vspace{8mm}
{\small\today\par}
\end{titlepage}

\tableofcontents
\newpage

% =====================================================================
\section{Overview}
\label{sec:overview}

\code{quizgen} takes an existing \textbf{question bank} --- a JSON list of
questions, each carrying metadata (chapter, difficulty, type, Bloom level) ---
and \emph{selects} from it the subset(s) that form a professor's exam. The
exam shape is described by a \textbf{blueprint}: how many questions in total,
how they split across chapters, the target difficulty mix, which question types
are allowed, and how many parallel, non-overlapping \emph{versions} to produce.

The scope is deliberately narrow: \textbf{there is no question generation and
no document ingestion}. The tool reasons only over the metadata of questions
it is given and the constraints it is asked to satisfy. This makes selection a
well-posed combinatorial optimisation problem with an exact, auditable
solution --- in contrast to prompting a language model to ``write the exam'',
which couples creation with constraint satisfaction and gives neither
guarantees nor traceability.

\paragraph{Why model selection as optimisation.}
Hitting a quantitative specification exactly --- ``10 from Chapter~1, 6 from
Chapter~2, 4 from Chapter~3; an 8/8/4 easy/medium/hard split; two versions that
share no question'' --- is a constraint-satisfaction-and-optimisation task.
Solving it with a Mixed-Integer Program gives:
\begin{itemize}[nosep,leftmargin=1.4em]
  \item \textbf{Exactness} --- the blueprint counts are met to the unit, or the
        run reports precisely which constraint had to be relaxed.
  \item \textbf{Multiple disjoint versions} --- produced in a single solve, so
        parallel exams provably share no items.
  \item \textbf{Quality-aware tie-breaking} --- among feasible selections, the
        one maximising a question-quality objective is chosen.
  \item \textbf{Reproducibility} --- the assembled exam embeds the resolved
        blueprint, so it is self-describing.
\end{itemize}

\paragraph{The stages.}
Figure~\ref{fig:pipeline} shows the flow. The pool is optionally
de-duplicated, assembled into one or more versions against the blueprint, and
each version is optionally validated. The stages map onto the modules in
Table~\ref{tab:modules}.

\begin{figure}[h]
\centering
\begin{tikzpicture}[node distance=10mm and 16mm]
  \node[data] (pool) {question bank\\(JSON)};
  \node[stage, right=of pool] (dedup) {Deduplicate\\(optional)};
  \node[stage, right=of dedup] (assemble) {Assemble\\(greedy / MIP)};
  \node[stage, right=of assemble] (validate) {Validate\\(optional)};
  \node[data, right=of validate] (out) {exam version(s)\\+ reports};
  \node[ctrl, above=10mm of assemble] (bp) {blueprint\\(YAML + flags)};

  \draw[arr] (pool) -- (dedup);
  \draw[arr] (dedup) -- (assemble);
  \draw[arr] (assemble) -- (validate);
  \draw[arr] (validate) -- (out);
  \draw[arr] (bp) -- (assemble);
\end{tikzpicture}
\caption{The selection pipeline. Green = data, blue = processing stage, red =
professor input. The blueprint drives assembly; dedup and validation are
opt-in.}
\label{fig:pipeline}
\end{figure}

\begin{table}[h]
\centering
\caption{Stage-to-module mapping.}
\label{tab:modules}
\small
\begin{tabular}{@{}llp{7.4cm}@{}}
\toprule
\textbf{Stage} & \textbf{Module} & \textbf{Responsibility} \\
\midrule
Data model    & \mod{schema.py}    & Pydantic models (\code{Question}, \code{Quiz}); JSON-Schema export. \\
Blueprint     & \mod{blueprint.py} & Blueprint model; count resolution (largest-remainder). \\
Deduplicate   & \mod{dedup.py}, \mod{similarity.py} & Near-duplicate detection (Jaccard / embedding); pool dedup. \\
Assemble      & \mod{assemble.py}  & Automated Test Assembly: greedy and MIP selectors. \\
Validate      & \mod{validate.py}, \mod{textmatch.py} & Schema, blueprint, grounding, duplicate checks. \\
CLI           & \mod{\_\_main\_\_.py} & Pool + blueprint/flags $\to$ exam version(s). \\
Config        & \mod{config.py}    & Single optional embedding-model setting. \\
\bottomrule
\end{tabular}
\end{table}

% =====================================================================
\section{Data model and contracts}
\label{sec:schema}

Every artifact is a Pydantic~v2 model (\mod{schema.py}), so the data contract
is the single source of truth and the JSON~Schema is \emph{exported} from it
(committed under \code{schemas/}) rather than hand-written --- the two cannot
drift. The central type is \code{Question}:

\begin{lstlisting}
class Question(BaseModel):
    id: str                 # uuid4 hex[:12]
    chapter: int            # >= 1
    section: str | None
    qtype: QuestionType     # mcq | true_false | short_answer | cloze
    bloom_level: BloomLevel # remember .. create  (revised Bloom, 2001)
    difficulty: Difficulty  # easy | medium | hard
    stem: str               # >= 10 chars
    options: list[str] | None
    answer: str             # >= 1 char
    explanation: str        # >= 10 chars
    source_ref: str         # provenance (free-form)
\end{lstlisting}

The metadata fields are exactly what selection reasons over
(Table~\ref{tab:fields}). The three enumerations make a tag impossible to
mistype. A \code{Quiz} wraps a list of questions plus the \code{blueprint} that
produced it, so an assembled exam is self-describing.

\begin{table}[h]
\centering
\caption{How each metadata field drives selection.}
\label{tab:fields}
\small
\begin{tabular}{@{}ll@{}}
\toprule
\textbf{Field} & \textbf{Used for} \\
\midrule
\code{chapter}                & per-chapter quotas \\
\code{difficulty}             & the easy / medium / hard mix \\
\code{qtype}                  & allowed-type filtering \\
\code{bloom\_level}           & the quality objective (favours discriminating levels) \\
\code{stem}/\code{answer}/\code{options} & near-duplicate detection \\
\code{source\_ref}            & provenance / audit trail \\
\bottomrule
\end{tabular}
\end{table}

% =====================================================================
\section{The blueprint}
\label{sec:blueprint}

A \emph{blueprint} (\mod{blueprint.py}) is the small YAML a professor edits ---
a \emph{table of specifications} for the exam:

\begin{lstlisting}[language=]
title: "AI-1 Midterm Exam"
total_questions: 20
chapters: { 1: "50%", 2: "30%", 3: "20%" }
difficulty_mix: { easy: 0.4, medium: 0.4, hard: 0.2 }
allowed_qtypes: [mcq, true_false, short_answer, cloze]
versions: 2
selector: mip            # greedy | mip
dedup_similarity: 0.85
\end{lstlisting}

\paragraph{Exact integer resolution.}
Chapter and difficulty amounts may be integer counts, fractions, or
\code{"NN\%"} strings. \code{Blueprint.resolve()} converts them to exact
integer counts that sum to \code{total\_questions} using the
\textbf{largest-remainder (Hamilton)} method:
\begin{equation}
\text{raw}_k = w_k \cdot N,\quad
\text{floor}_k = \lfloor \text{raw}_k \rfloor,\quad
r = N - \!\sum_k \text{floor}_k ,
\end{equation}
after which the $r$ leftover units go to the $k$ with the largest fractional
remainders $\text{raw}_k - \text{floor}_k$. This guarantees, for instance, that
``40/40/20 of 20'' resolves unambiguously to $8/8/4$.

\paragraph{YAML and flags (flags win).}
The same constraints can be supplied or overridden on the command line, so a
base \code{blueprint.yaml} can be tweaked per run. \code{build\_blueprint} loads
the YAML (if any) into a dict and overwrites field-by-field with any provided
flags before validating the \code{Blueprint}:

\begin{lstlisting}[language=bash]
python -m quizgen --pool questions.json --blueprint blueprint.yaml \
       --versions 3 --selector greedy        # override just these two
\end{lstlisting}

% =====================================================================
\section{Automated Test Assembly}
\label{sec:assemble}

Selection is \emph{Automated Test Assembly} (ATA): choose, from the pool, the
subset(s) satisfying all blueprint constraints while maximising overall
quality. Two interchangeable selectors sit behind one \code{Selector} interface
(chosen by \code{blueprint.selector}):

\begin{description}[leftmargin=1.4em,style=nextline]
  \item[\code{GreedySelector}] A fast, content-balancing greedy fill. Honours
  per-chapter counts, respects the difficulty marginal in a two-pass fill
  (quota first, then remainder), and skips near-duplicates and questions already
  used in another version. A good baseline with no optimality guarantee.
  \item[\code{MIPSelector}] A Mixed-Integer Program (PuLP\,/\,CBC) that meets
  per-chapter \emph{and} difficulty counts exactly, maximises total quality,
  keeps parallel versions disjoint, and forbids near-duplicate pairs across
  versions. This is the exact, auditable selection.
\end{description}

\paragraph{MIP formulation.}
Let $\mathcal{P}$ be the (qtype-filtered) pool, $V$ the set of versions, and
$x_{iv}\in\{0,1\}$ indicate that question $i$ is placed in version $v$. With
per-chapter targets $n_c$, per-difficulty targets $m_d$, a quality weight
$q_i$, and the set $D$ of near-duplicate pairs:

\begin{align}
\max \quad & \sum_{v\in V}\sum_{i\in\mathcal{P}} q_i\, x_{iv}
   &&\text{(maximise total quality)} \label{eq:obj}\\
\text{s.t.}\quad
& \sum_{i:\,\text{chap}(i)=c} x_{iv} = n_c
   && \forall v\in V,\ \forall c \label{eq:chap}\\
& \sum_{i:\,\text{diff}(i)=d} x_{iv} = m_d
   && \forall v\in V,\ \forall d \label{eq:diff}\\
& \sum_{v\in V} x_{iv} \le 1
   && \forall i\in\mathcal{P} \label{eq:disjoint}\\
& \sum_{v\in V} x_{iv} + \sum_{v\in V} x_{jv} \le 1
   && \forall (i,j)\in D \label{eq:dup}
\end{align}

Constraint~\eqref{eq:chap} fixes the per-chapter blueprint counts;
\eqref{eq:diff} fixes the difficulty mix; \eqref{eq:disjoint} makes the
parallel versions \emph{disjoint} (no question reused); \eqref{eq:dup} forbids
any near-duplicate \emph{pair} from appearing anywhere across the assembled
set. The objective~\eqref{eq:obj} ranks otherwise-feasible selections by a
cheap quality proxy (\code{question\_quality}: rewards substantive
explanations, well-formed 4-option MCQs, and the more discriminating Bloom
levels).

\paragraph{Graceful infeasibility.}
If the exact program is infeasible for the given pool (e.g.\ too few
\code{hard} questions in some chapter), the solver \emph{relaxes} the difficulty
equalities~\eqref{eq:diff}, records \code{difficulty\_relaxed} in the
diagnostics, and re-solves --- the per-chapter counts and totals are never
silently violated. Capacity warnings are emitted up-front when the pool is too
small to honour the demand.

\begin{figure}[h]
\centering
\begin{tikzpicture}[node distance=7mm and 14mm]
  \node[data] (pool) {tagged pool\\($|\mathcal{P}|$ questions)};
  \node[ctrl, right=of pool] (bp) {blueprint\\$n_c,\,m_d,\,|V|$};
  \coordinate (mid) at ($(pool)!0.5!(bp)$);
  \node[stage, below=10mm of mid, minimum width=40mm] (mip)
        {MIP: maximise $\sum q_i x_{iv}$\\s.t.\ \eqref{eq:chap}--\eqref{eq:dup}};
  \node[data, below left=10mm and -6mm of mip] (vA) {Version A};
  \node[data, below right=10mm and -6mm of mip] (vB) {Version B};

  \draw[arr] (pool) -- (mip);
  \draw[arr] (bp) -- (mip);
  \draw[arr] (mip) -- (vA);
  \draw[arr] (mip) -- (vB);
  \node[font=\tiny,below=1mm of vA,align=center] {exact counts,\\disjoint, no dups};
\end{tikzpicture}
\caption{MIP assembly: one optimisation produces all parallel versions at once,
guaranteeing they are disjoint and individually blueprint-compliant.}
\label{fig:mip}
\end{figure}

% =====================================================================
\section{Deduplication}
\label{sec:dedup}

Near-duplicate questions waste blueprint slots and make parallel versions
predictable. \mod{similarity.py} provides one API
(\code{find\_near\_duplicate\_pairs}) with two backends:

\begin{itemize}[nosep,leftmargin=1.4em]
  \item \textbf{Token-Jaccard} over the canonical question text
  (stem~+~answer~+~MCQ options) --- dependency-free, always available:
  \[
    \mathrm{sim}(q_1,q_2) = \frac{|T(q_1)\cap T(q_2)|}{|T(q_1)\cup T(q_2)|},
  \]
  where $T(\cdot)$ is the set of lowercased word tokens.
  \item \textbf{Embedding cosine} using an optional sentence-transformers
  encoder, which catches paraphrases the lexical measure misses. Cosine
  similarity is computed on $L_2$-normalised vectors.
\end{itemize}

Pairs scoring at/above a threshold are returned as index pairs. The optional
pool-dedup pass (\mod{dedup.py}, CLI \code{--dedup}) drops the later member of
each pair \emph{before} assembly; independently, the assembler forbids any
surviving near-duplicate pair across versions (constraint~\eqref{eq:dup}). The
backend degrades gracefully: with no encoder available it uses the offline
token measure automatically.

% =====================================================================
\section{Validation}
\label{sec:validate}

\code{validate\_exam} (\mod{validate.py}, CLI \code{--validate}) produces a
structured report over four dimensions:

\begin{enumerate}[nosep,leftmargin=1.6em]
  \item \textbf{Schema validity} --- every question round-trips through the
        \code{Question} model.
  \item \textbf{Blueprint compliance} --- per-chapter counts, difficulty mix,
        total, and allowed qtypes match the blueprint \emph{exactly}.
  \item \textbf{Grounding coverage} --- share of questions carrying a
        \code{source\_ref} (and, when source text is supplied, whose answer
        words overlap the referenced passage).
  \item \textbf{Duplicate rate} --- share of questions involved in a
        near-duplicate pair.
\end{enumerate}

The overall gate \code{passed} is true only if no schema-invalid items exist
and (if checked) the blueprint is satisfied. The grounding check reuses the
dependency-free \code{extract\_content\_words} helper in \mod{textmatch.py},
kept separate so validation pulls in no heavy dependencies.

% =====================================================================
\section{Command-line interface}
\label{sec:cli}

\mod{\_\_main\_\_.py} is the single entry point. It loads the pool, builds the
blueprint from the YAML and/or flag overrides, optionally de-duplicates,
assembles, writes one \code{Quiz} JSON per version (each embedding the resolved
blueprint), and --- with \code{--validate} --- writes a Markdown report per
version. It exits non-zero if any version is non-compliant, so it composes in
scripts and CI.

\begin{lstlisting}[language=bash]
# Pure command line: 20 questions, 10/6/4 across chapters, 40/40/20
# difficulty, two disjoint versions, optimal (MIP) selection:
python -m quizgen --pool questions.json \
    --total 20 --chapters 1:10,2:6,3:4 \
    --difficulty 0.4,0.4,0.2 --versions 2 --selector mip

# Blueprint file + dedup + per-version validation report:
python -m quizgen --pool questions.json --blueprint blueprint.yaml \
    --dedup --validate --out-dir out
\end{lstlisting}

The assembly and validation steps also expose their own entry points
(\code{python -m quizgen.assemble}, \code{python -m quizgen.validate}) for use
as standalone tools.

% =====================================================================
\section{Provenance of techniques}
\label{sec:sources}

The table below states, for each implemented technique, the body of literature
it draws on. \emph{Honesty note:} the citation keys are the standard references
for each method and \textbf{should be replaced with the exact bibliography
entries the thesis adopts}. The right-hand column distinguishes ideas from the
literature from project-specific engineering.

\begin{longtable}{@{}p{4.1cm}p{4.3cm}p{4.4cm}@{}}
\toprule
\textbf{Implemented in code} & \textbf{Technique / idea} & \textbf{Source lineage} \\
\midrule
\endhead
\mod{assemble.py} (\code{MIPSelector}) &
Automated Test Assembly as 0--1 MIP &
van der Linden, \emph{Linear Models for Optimal Test Design} (2005); exam
assembly as integer programming. \\[2pt]
\code{question\_quality} objective &
Objective-weighted selection in ATA &
van der Linden (2005); here a lightweight quality proxy rather than IRT
information. \\[2pt]
\mod{blueprint.py} &
Test blueprint / table of specifications &
Classical assessment design; content $\times$ cognitive-level specification
driving item selection. \\[2pt]
\code{BloomLevel} metadata &
Revised Bloom's taxonomy &
Anderson \& Krathwohl (2001), revising Bloom (1956); cognitive-level tag used
in the quality objective. \\[2pt]
\mod{similarity.py} (embedding cosine) &
Semantic near-duplicate detection &
Sentence-BERT embeddings (Reimers \& Gurevych, 2019). \\
\midrule
\multicolumn{3}{@{}l}{\emph{Project-specific engineering (not from a paper):}}\\
\midrule
\code{\_largest\_remainder} &
Hamilton / largest-remainder apportionment &
Standard apportionment arithmetic; ensures integer counts sum exactly. \\[2pt]
\code{GreedySelector} &
Greedy content-balanced baseline &
Engineering baseline to contrast against the MIP. \\[2pt]
Jaccard fallback, \mod{textmatch.py} &
Lexical token overlap &
Dependency-free fallbacks for offline operation. \\[2pt]
YAML$+$flag blueprint, CLI &
Configuration / override precedence &
Standard tooling; flags override file defaults. \\
\bottomrule
\end{longtable}

% =====================================================================
\section{Summary}

\code{quizgen} implements the \emph{selection} half of automated exam
construction as an exact optimisation problem. Given a tagged question bank and
a professor's blueprint, it assembles blueprint-compliant, mutually disjoint
exam versions via a 0--1 Mixed-Integer Program, with a greedy baseline behind
the same interface; an optional deduplication pass and a layered validation
report (schema, blueprint, grounding, duplicates) verify the result. The tool
has no LLM or document-processing dependencies, runs fully offline, and exposes
its constraints through a blueprint file and/or command-line flags.

\end{document}
